\section{Определение скалярного произведения в унитарном пространстве}

\begin{definition}[Скалярное произведение в унитарном пространстве]
Пусть $V$ --- линейное пространство над $\mathbb{C}$.

$V$ называется \textbf{унитарным}, если в нём введена функция $(x, y): V \times V \to \mathbb{C}$ (\textbf{скалярное произведение}), удовлетворяющая условиям:
\begin{enumerate}
    \item $(x, y) = \overline{(y, x)},\quad \forall x, y \in V$;
    \item $(x_1 + x_2, y) = (x_1, y) + (x_2, y),\quad \forall x_1, x_2, y \in V$;
    \item $(\lambda x, y) = \lambda (x, y),\quad \forall \lambda \in \mathbb{C},\ \forall x, y \in V$;
    \item $(x, x) \in \mathbb{R}$, $(x, x) > 0$ при $x \neq 0$; $(x, x) = 0 \iff x = 0$.
\end{enumerate}
\end{definition}
